% GENERATED FILE. Edit canonical lessons, then run npm run generate:books.
% canonical-source-hash: fnv1a64:5410a2d994c5fc66
% canonical-lessons: PA-C13-dil, PA-C13-sir

\chapter{Heart and Head}
\label{ch:pa-heart-and-head}
\hlchaptermodality{13}

\begin{chapteropening}
\textbf{By the end of this chapter:} \emph{I can name the heart and the head in Punjabi, tell that \pa{ਦਿਲ} is Punjabi's one borrowed body word against four inherited ones, and trace \pa{ਸਿਰ} to the same root as English horn rather than English head.}

The tranche closes on \pa{ਸਿਰ}: the reader produces all six body words from this tranche, states that \pa{ਸਿਰ}'s real English cousin is horn, not head, and that its resemblance to Persian sar is convergence from a shared Proto-Indo-Iranian ancestor rather than a loan — the opposite case from \pa{ਦਿਲ}, which genuinely was borrowed from Persian despite sharing the same deep Indo-European root as English heart.
\end{chapteropening}

\section[dil]{\pa{ਦਿਲ} — the one body word that is not Sanskrit}

\label{lesson:PA-C13-dil}



\begin{quote}
\textbf{\pa{ਅੱਖ}}, \textbf{\pa{ਕੰਨ}}, \textbf{\pa{ਮੂੰਹ}}, \textbf{\pa{ਨੱਕ}} — four body words, all inherited from Sanskrit. This one breaks the pattern the same way \textbf{\pa{ਦੋਸਤ}} once did.
\end{quote}

\subsection*{You'll want to know first — one word}
\begin{quote}
\textbf{\pa{ਦਿਲ}} — \emph{dil} — \textbf{heart}
\end{quote}

\subsection*{The letters in this word}
\textbf{\pa{ਦ}} \emph{d}, \textbf{\pa{ਿ}} short \emph{i}, \textbf{\pa{ਲ}} \emph{l}. Every letter already yours, from \textbf{\pa{ਦੁੱਧ}} and \textbf{\pa{ਲੈਣਾ}}.

\begin{cousinweb}
\textbf{\pa{ਦਿਲ}} is \textbf{Persian}, borrowed the way \textbf{\pa{ਦੋਸਤ}} was — Punjabi's everyday word for \textquotedblleft{}heart\textquotedblright{} is not the word Sanskrit handed down. Sanskrit's own word, \textbf{hṛdaya}, exists in Punjabi too, but only as a formal, literary borrowing; ordinary speech reaches for \emph{dil}.

And yet \emph{dil} is not a stranger to Sanskrit's word. Both descend from the same Proto-Indo-European root, *\emph{kerd-}, \textquotedblleft{}heart\textquotedblright{} — Persian's branch carried it to \emph{dil}, Sanskrit's branch carried it to \emph{hṛdaya}, and a third branch carried it, unbroken, to English \textbf{heart} itself. Three languages, three shapes, one ancestor. What makes \emph{dil} worth flagging is not the deep root — that root is shared honestly — but the fact that Punjabi took its \emph{everyday} word for heart by borrowing sideways from Persian, rather than by inheriting straight down from Sanskrit the way \emph{akkh} and \emph{nakk} did.
\end{cousinweb}

\subsection*{Your turn}
\begin{itemize}
\raggedright
  \item \emph{Say it:} \textbf{dil} — heart
  \item \emph{Trace it:} PIE \emph{kerd-} $\to$ Persian \emph{dil} / Sanskrit \emph{hṛdaya} / English \textbf{heart}, three branches
  \item \emph{Contrast:} \emph{dil}, borrowed, against \emph{akkh} and \emph{nakk}, inherited
  \item \emph{Recall:} the other Persian loan this book has taught for a person you know (\textbf{dost})
\end{itemize}

\begin{tcolorbox}[breakable,colback=teal!4,colframe=teal!35!black,title={Before you move on}]
Is \emph{dil} inherited from Sanskrit or borrowed from Persian? (\textbf{Borrowed}, from Persian — Sanskrit's own word, \emph{hṛdaya}, stays formal.) What single Proto-Indo-European root sits behind \emph{dil}, Sanskrit \emph{hṛdaya}, and English \emph{heart}, all three? (*\emph{Kerd-}.)

Sources: \href{https://en.wiktionary.org/wiki/\%E0\%A8\%A6\%E0\%A8\%BF\%E0\%A8\%B2}{Wiktionary: \pa{ਦਿਲ}}; \href{https://en.wiktionary.org/wiki/\%D8\%AF\%D9\%84}{Wiktionary: Persian dil}.
\end{tcolorbox}
\section[sir]{\pa{ਸਿਰ} — a cousin of \textquotedblleft{}horn,\textquotedblright{} and a twin nobody borrowed}

\label{lesson:PA-C13-sir}



\begin{quote}
\textbf{\pa{ਦਿਲ}} broke the pattern — borrowed, where every other body word was inherited. This last word restores the pattern, and closes the tranche on its most surprising cousin yet. On the way, say \textbf{\pa{ਭਰਾ}} and \textbf{\pa{ਭੈਣ}} once more, low-toned, the way this book's other loan-versus-inheritance pair opened it.
\end{quote}

\subsection*{You'll want to know first — one word}
\begin{quote}
\textbf{\pa{ਸਿਰ}} — \emph{sir} — \textbf{head}
\end{quote}

\subsection*{The letters in this word}
\textbf{\pa{ਸ}} \emph{s}, \textbf{\pa{ਿ}} short \emph{i}, \textbf{\pa{ਰ}} \emph{r}. Every letter already familiar, from \textbf{\pa{ਸਤਿ}} and \textbf{\pa{ਮੇਰਾ}}.

\begin{cousinweb}
\textbf{\pa{ਸਿਰ}} is inherited, back to Proto-Indo-European *\emph{kerh\textsubscript{2}-}, \textquotedblleft{}head, horn.\textquotedblright{} That is genuinely surprising: English \textbf{head} comes from an entirely different, unrelated root, *\emph{kaput-}. \emph{Sir}'s real English cousin is \textbf{horn} — the same *\emph{kerh\textsubscript{2}-} that gave Latin \emph{cornū} and the Greek word for \textquotedblleft{}head,\textquotedblright{} \emph{karēnon}. A head, on this root's evidence, was once named for whatever stood up highest and hardest on it.

One more thing \emph{sir} does, that no word in this book has done before: it looks identical to Persian \textbf{sar}, \textquotedblleft{}head\textquotedblright{} — and unlike \emph{dil}, that is \textbf{not} a loan. Persian and Sanskrit are sister branches of Proto-Indo-Iranian, and each inherited *\emph{kerh\textsubscript{2}-} on its own, straight down its own line. \emph{Dil} is Punjabi reaching sideways into Persian for a word it did not otherwise have. \emph{Sir} only looks like the same kind of borrowing; it is a twin by shared birth, not by loan — closing this tranche the same way Chapter 6 closed on \emph{panj}, two branches independently arriving at one shape.

Running the whole face and the two words beyond it: \textbf{akkh}, \textbf{kann}, \textbf{mūnh}, \textbf{nakk} are all inherited and tone-marked only where a consonant was lost outright; \textbf{dil} is the one loanword; \textbf{sir} is inherited, and its resemblance to Persian is convergence, exactly like \emph{panj}'s.
\end{cousinweb}

\subsection*{Your turn}
\begin{itemize}
\raggedright
  \item \emph{Say it:} \textbf{sir} — head
  \item \emph{State:} \emph{sir}'s real English cousin (\textbf{horn}, not \textbf{head})
  \item \emph{Contrast:} \emph{sir} and Persian \emph{sar} — look-alikes by shared ancestry, not by loan; \emph{dil} — a genuine loan
  \item \emph{Run through:} all six body words from this tranche — \textbf{akkh}, \textbf{kann}, \textbf{mūnh}, \textbf{nakk}, \textbf{dil}, \textbf{sir}
\end{itemize}

\begin{tcolorbox}[breakable,colback=teal!4,colframe=teal!35!black,title={Before you move on}]
Give the word for \textquotedblleft{}head.\textquotedblright{} (\emph{Sir}.) Which English word is its real cousin — \emph{head}, or \emph{horn}? (\textbf{Horn}.) Is \emph{sir}'s resemblance to Persian \emph{sar} a loan, the way \emph{dil} is? (\textbf{No} — both inherited the same root separately, the way \emph{panj} and Persian \emph{panj} once did.) Name all six body words from this tranche. (\emph{Akkh}, \emph{kann}, \emph{mūnh}, \emph{nakk}, \emph{dil}, \emph{sir}.)

Sources: \href{https://en.wiktionary.org/wiki/\%E0\%A8\%B8\%E0\%A8\%BF\%E0\%A8\%B0}{Wiktionary: \pa{ਸਿਰ}}; \href{https://en.wiktionary.org/wiki/Reconstruction:Proto-Indo-European/\%E1\%B8\%B1erh\%E2\%82\%82-}{Wiktionary: \'{k}erh\textsubscript{2}-}.
\end{tcolorbox}
